% GENERATED FILE. Edit canonical lessons, then run npm run generate:books.
% canonical-source-hash: fnv1a64:700c3f10fd53a728
% canonical-lessons: HI-C293-thagna, HI-C293-dhokha, HI-C293-risvat, HI-C293-badmas, HI-C293-naqli

\chapter{Cheats}
\label{ch:hi-cheats}
\hlchaptermodality{293}

\begin{chapteropening}
\textbf{By the end of this chapter:} \emph{I can say how someone was cheated and what to do if it happens.}

Complete the last lesson of chapter 293: I can say how someone was cheated and what to do if it happens.
\end{chapteropening}

\section[\texorpdfstring{\hi{ठगना} (ṭhagnā)}{ठगना (ṭhagnā)}]{\hi{ठगना} (ṭhagnā) — to cheat; to swindle}

\label{lesson:HI-C293-thagna}



\begin{quote}
Before the new one: say the Hindi for a gun, then the Hindi for to tie up.
\end{quote}

\subsection*{You'll want to know: \hi{ठगना}}
\textbf{\hi{ठगना}} — \emph{ṭhagnā} — \textquotedblleft{}to cheat; to swindle\textquotedblright{}.

\begin{quote}
\hi{अगर} \hi{कोई} \hi{तुम्हें} \hi{ठग} \hi{रहा} \hi{है}, \hi{तो} \hi{पुलिस} \hi{को} \hi{बताओ।}
\end{quote}

\begin{quote}
\emph{agar koī tumhẽ ṭhag rahā hai, to pulis ko batāo.}
\end{quote}

\begin{quote}
If someone is cheating you, tell the police.
\end{quote}

Past, with \hi{ने} (he or she …): \hi{उसने} \hi{मुझे} \hi{ठगा}, \emph{usne mujhe ṭhagā}.

\begin{cousinweb}
\textbf{\hi{ठगना}} (\emph{ṭhagnā}) is inherited from \emph{*ṭhagg}, a form CDIAL reconstructs (CDIAL 5489), which also lists the noun \emph{ṭhag}, a swindler. English \emph{thug} was taken from that noun.
\end{cousinweb}

\begin{grammarlens}[title={What you've built}]
One more word for telling what happened.
\end{grammarlens}

\subsection*{Your turn}
\begin{itemize}
\raggedright
  \item \emph{Say it:} \emph{ṭhagnā}
  \item \emph{Say it:} \emph{ṭhagnā}, once more
  \item \emph{Say it:} \emph{agar koī tumhẽ ṭhag rahā hai, to pulis ko batāo.} — If someone is cheating you, tell the police.
  \item \emph{Recall:} say the Hindi for a gun, then the Hindi for to tie up, then say \emph{ṭhagnā} again
\end{itemize}

\begin{tcolorbox}[breakable,colback=teal!4,colframe=teal!35!black,title={Before you move on}]
What does \hi{ठगना} mean? (\textquotedblleft{}To cheat; to swindle\textquotedblright{}.) Say it once more.
\end{tcolorbox}
\section[\texorpdfstring{\hi{धोखा} (dhokhā)}{धोखा (dhokhā)}]{\hi{धोखा} (dhokhā) — a deception; a trick}

\label{lesson:HI-C293-dhokha}



\begin{quote}
Before the new one: say the Hindi for to tie up, then the Hindi for to cheat.
\end{quote}

\subsection*{You'll want to know: \hi{धोखा}}
\textbf{\hi{धोखा}} — \emph{dhokhā} — \textquotedblleft{}a deception; a trick\textquotedblright{}.

\begin{quote}
\hi{उस} \hi{आदमी} \hi{ने} \hi{हमें} \hi{धोखा} \hi{दिया।}
\end{quote}

\begin{quote}
\emph{us ādmī ne hamẽ dhokhā diyā.}
\end{quote}

\begin{quote}
That man deceived us.
\end{quote}

Masculine.

\begin{cousinweb}
\textbf{\hi{धोखा}} (\emph{dhokhā}) is inherited from \emph{*dhrōkṣa}, a form CDIAL reconstructs from its descendants (CDIAL 6894). \textbf{\hi{धोखा} \hi{देना}} is to deceive someone; \textbf{\hi{धोखा} \hi{खाना}}, to be taken in.
\end{cousinweb}

\begin{grammarlens}[title={What you've built}]
One more word for telling what happened.
\end{grammarlens}

\subsection*{Your turn}
\begin{itemize}
\raggedright
  \item \emph{Say it:} \emph{dhokhā}
  \item \emph{Say it:} \emph{dhokhā}, once more
  \item \emph{Say it:} \emph{us ādmī ne hamẽ dhokhā diyā.} — That man deceived us.
  \item \emph{Recall:} say the Hindi for to tie up, then the Hindi for to cheat, then say \emph{dhokhā} again
\end{itemize}

\begin{tcolorbox}[breakable,colback=teal!4,colframe=teal!35!black,title={Before you move on}]
What does \hi{धोखा} mean? (\textquotedblleft{}A deception; a trick\textquotedblright{}.) Say it once more.
\end{tcolorbox}
\section[\texorpdfstring{\hi{रिश्वत} (riśvat)}{रिश्वत (riśvat)}]{\hi{रिश्वत} (riśvat) — a bribe}

\label{lesson:HI-C293-risvat}



\begin{quote}
Before the new one: say the Hindi for to cheat, then the Hindi for a deception.
\end{quote}

\subsection*{You'll want to know: \hi{रिश्वत}}
\textbf{\hi{रिश्वत}} — \emph{riśvat} — \textquotedblleft{}a bribe\textquotedblright{}.

\begin{quote}
\hi{अगर} \hi{वह} \hi{रिश्वत} \hi{माँगेगा}, \hi{तो} \hi{मैं} \hi{पुलिस} \hi{को} \hi{बताऊँगा।}
\end{quote}

\begin{quote}
\emph{agar vah riśvat mā̃gegā, to maiṁ pulis ko batāūṅgā.}
\end{quote}

\begin{quote}
If he asks for a bribe, I will tell the police.
\end{quote}

Feminine.

\begin{cousinweb}
\textbf{\hi{रिश्वत}} (\emph{riśvat}) is usually traced, through Persian, to Arabic \emph{rishwa}, \textquotedblleft{}a bribe\textquotedblright{}. It is feminine: \textbf{\hi{उसने} \hi{रिश्वत} \hi{माँगी}}, he asked for a bribe.
\end{cousinweb}

\begin{grammarlens}[title={What you've built}]
One more word for telling what happened.
\end{grammarlens}

\subsection*{Your turn}
\begin{itemize}
\raggedright
  \item \emph{Say it:} \emph{riśvat}
  \item \emph{Say it:} \emph{riśvat}, once more
  \item \emph{Say it:} \emph{agar vah riśvat mā̃gegā, to maiṁ pulis ko batāūṅgā.} — If he asks for a bribe, I will tell the police.
  \item \emph{Recall:} say the Hindi for to cheat, then the Hindi for a deception, then say \emph{riśvat} again
\end{itemize}

\begin{tcolorbox}[breakable,colback=teal!4,colframe=teal!35!black,title={Before you move on}]
What does \hi{रिश्वत} mean? (\textquotedblleft{}A bribe\textquotedblright{}.) Say it once more.
\end{tcolorbox}
\section[\texorpdfstring{\hi{बदमाश} (badmāś)}{बदमाश (badmāś)}]{\hi{बदमाश} (badmāś) — a crook; a rascal}

\label{lesson:HI-C293-badmas}



\begin{quote}
Before the new one: say the Hindi for a deception, then the Hindi for a bribe.
\end{quote}

\subsection*{You'll want to know: \hi{बदमाश}}
\textbf{\hi{बदमाश}} — \emph{badmāś} — \textquotedblleft{}a crook; a rascal\textquotedblright{}.

\begin{quote}
\hi{दो} \hi{बदमाश} \hi{बस} \hi{में} \hi{लोगों} \hi{का} \hi{पैसा} \hi{छीन} \hi{रहे} \hi{थे।}
\end{quote}

\begin{quote}
\emph{do badmāś bas meṁ logoṁ kā paisā chīn rahe the.}
\end{quote}

\begin{quote}
Two crooks were snatching people's money on the bus.
\end{quote}

Masculine or feminine, by the person meant.

\begin{cousinweb}
\textbf{\hi{बदमाश}} (\emph{badmāś}) is usually traced to Persian \emph{bad}, \textquotedblleft{}bad\textquotedblright{}, and \emph{maʿāš}, \textquotedblleft{}a livelihood\textquotedblright{}, a word Persian took from Arabic: someone who lives by bad means. Said fondly to a child, it means a little rascal.
\end{cousinweb}

\begin{grammarlens}[title={What you've built}]
One more word for telling what happened.
\end{grammarlens}

\subsection*{Your turn}
\begin{itemize}
\raggedright
  \item \emph{Say it:} \emph{badmāś}
  \item \emph{Say it:} \emph{badmāś}, once more
  \item \emph{Say it:} \emph{do badmāś bas meṁ logoṁ kā paisā chīn rahe the.} — Two crooks were snatching people's money on the bus.
  \item \emph{Recall:} say the Hindi for a deception, then the Hindi for a bribe, then say \emph{badmāś} again
\end{itemize}

\begin{tcolorbox}[breakable,colback=teal!4,colframe=teal!35!black,title={Before you move on}]
What does \hi{बदमाश} mean? (\textquotedblleft{}A crook; a rascal\textquotedblright{}.) Say it once more.
\end{tcolorbox}
\section[\texorpdfstring{\hi{नक़ली} (naqlī)}{नक़ली (naqlī)}]{\hi{नक़ली} (naqlī) — fake; counterfeit}

\label{lesson:HI-C293-naqli}



\begin{quote}
Before the new one: say the Hindi for a bribe, then the Hindi for a crook.
\end{quote}

\subsection*{You'll want to know: \hi{नक़ली}}
\textbf{\hi{नक़ली}} — \emph{naqlī} — \textquotedblleft{}fake; counterfeit\textquotedblright{}.

\begin{quote}
\hi{उसने} \hi{मुझे} \hi{एक} \hi{नक़ली} \hi{नोट} \hi{दिया।}
\end{quote}

\begin{quote}
\emph{usne mujhe ek naqlī noṭ diyā.}
\end{quote}

\begin{quote}
He gave me a fake banknote.
\end{quote}

\begin{cousinweb}
\textbf{\hi{नक़ली}} (\emph{naqlī}) is \emph{naql}, a copy, usually traced through Persian to Arabic \emph{naql}, \textquotedblleft{}a copying, a carrying over\textquotedblright{}, with the adjective ending \emph{-ī}: a copy passed off as the real thing.
\end{cousinweb}

\begin{grammarlens}[title={What you've built}]
One more word for telling what happened.
\end{grammarlens}

\subsection*{Your turn}
\begin{itemize}
\raggedright
  \item \emph{Say it:} \emph{naqlī}
  \item \emph{Say it:} \emph{naqlī}, once more
  \item \emph{Say it:} \emph{usne mujhe ek naqlī noṭ diyā.} — He gave me a fake banknote.
  \item \emph{Recall:} say the Hindi for a bribe, then the Hindi for a crook, then say \emph{naqlī} again
\end{itemize}

\begin{tcolorbox}[breakable,colback=teal!4,colframe=teal!35!black,title={Before you move on}]
What does \hi{नक़ली} mean? (\textquotedblleft{}Fake; counterfeit\textquotedblright{}.) Say it once more.
\end{tcolorbox}
